\scanpage{11}\mspage{10}
Pour ceci, il \inserted{\unsure{faut et|f. et}} suffit qu'on ait l'axiome suivant

% slanted
\margin{Critère de compatibilité absolue des ateliers spacieux}
\uline{At CD 2}\quad Soient $x,y\deleted{\unsure{,}}\in\unsure{M|\mathcal{M}}$, \deleted{\unsure{soit}} on ne suppose pas $x\deleted{\mathrel{>\mkern-12mu<}}y$ compatibles mais on introduit $K=x\cap y=\displaystyle\mathop{\mathrm{Sup}}_{\unsure{z|Z}\,\in\,\widetilde x\cap\widetilde y}z$. On suppose que pour tout $x',y'\in M$ avec $x'\ll x$, $x'_L=\unsure{\phi|\varnothing}_{\mathcal{F}}$ et $y'\ll y$, $y'_L=\phi_{\mathcal{F}}$, on ait $x'\deleted{\mathrel{>\mkern-12mu<}}\parallel y'$. \deleted{\inserted{\illegible{2}}} Alors $x\mathrel{>\mkern-12mu<}y$.

Cet axiome n'est pas satisfait \inserted{le plus \unsure{souvent}} pour les \uline{ateliers} dits ``\uline{\unsure{polyédraux}}'', i.e. satisfaisant \unsure{à :|:} \deleted{\unsure{pour} tout} \inserted{deux} multistrates, \deleted{\unsure{doivent être des} \illegible{1}} \inserted{compatibles, minorés} \deleted{\unsure{disjointes|disjoints}} ont pour intersection une \unsure{multistrate}
i.e. si $x,y$ \deleted{$\unsure{\in M}$} dans $M$ \deleted{et \illegible{1} $x,y$ \unsure{majorés}} \inserted{avec $x\mathrel{>\mkern-12mu<}y$}, \deleted{\unsure{et}} \inserted{et} minorés dans $M$ \inserted{(pas néc. \unsure{majorés})}, alors \deleted{inf dans $M$} $x,y$ ont un inf dans $M$ i.e. $x\underset{\mathcal{F}}{\cap}y\in M$. Pour ceux-ci, \unsure{posera}

% slanted
\margin{Critère de compatibilité absolue \unsure{des|aux|dans} ateliers \unsure{polyédraux}}
\uline{At CD} $2'$ \quad Soient $x,y\in M$, $K=x\cap y=\displaystyle\mathop{\mathrm{Sup}}_{z\,\in\,\widetilde x\cap\widetilde y}z$. On suppose que $\forall\,x',y'\in M$, $x'\ll x$, $x'_L=\phi_{\mathcal{F}}$, $y'\ll y$, $y'_L=\phi_{\mathcal{F}}$, on ait $x'\parallel y'$. \deleted{Alors $x\mathrel{>\mkern-12mu<}y$.} Supposons de plus $K\in M$ i.e. c'est un $\mathrm{inf}^{\unsure{\leq}}$ de $x,y$ dans $\deleted{\unsure{Z}}M$. Alors $x\mathrel{>\mkern-12mu<}y$.

\deleted{\uline{NB} Chacune de ces deux conditions implique \uline{At CD} $1_0$.}

\noindent\begin{tabular}{||p{.92\linewidth}}
\uline{N.B.} \unsure{Tous} les ateliers que j'ai \unsure{vus} comme \unsure{motivants} satisfont At CD $1_0$, et l'un des deux axiomes At CD 2 ou At. CD $2'$.
\end{tabular}

% slanted
\margin{Attention La relation $\parallel$ \uline{dans $M$} ne résulte pas (comme dans $\unsure{\overline{F}|\mathcal{F}}$) de la seule donnée de $\leq$ -- ici il faut la voir (\unsure{en|ou} l'absence \unsure{en général|en la définition} de $\unsure{\widehat F|\overline F}$ \unsure{ou|en} $\unsure{b_{\overline{0}}|b_0}$) comme une \unsure{donnée} \uline{primitive}.}
L'intérêt de At CD, c'est de ramener la relation $\mathrel{>\mkern-12mu<}$ dans $M$, essentielle pour la création des figures, à la propriété $\parallel$ de disjonction dans $M$, \deleted{qui s'exprime en termes de $\leq$} \unsure{tel est que|tel est} le rôle crucial de cette notion de disjonction dans la construction des ateliers : partir de $(M,\leq,\ll\deleted{\unsure{,\mathrel{>\mkern-12mu<}}})$, \deleted{\unsure{seulement,}} \unsure{ou|on} maintenant en termes de $(M,\leq,\ll,\parallel_M)$
